\input{./._relpath-to-latexroot.ltx}
\documentclass[\latexroot/\projectname]{subfiles}
\input{\latexroot/@resources/subfile-setup.ltx}

\begin{document}

\section{Model}\whenintegrated{\label{sec:model}}

A cost-minimising government chooses how to allocate a given stock of
debt across three creditor sectors --- domestic non-banks, foreigners,
and the domestic banking system (central bank plus private banks,
treated as one block). The three sectors differ in \emph{how}
placement becomes costly at the margin, and those differences produce
a tractable pecking order with sharp comparative-static predictions
that we take to the data.

%-------------------------------------------------------------
\subsection{Setting}
%-------------------------------------------------------------

The government has a given stock of outstanding debt $D>0$ and chooses
non-negative holdings $(d_h,d_f,d_b)$ satisfying the budget constraint
\begin{equation}\label{eq:budget}
  d_h + d_f + d_b = D,
\end{equation}
where $d_h$ is non-bank holdings, $d_f$ is foreign holdings, and $d_b$
is banking absorption. Each sector is summarised by the cost of
placing debt with it. Non-banks and foreigners lend at interest rates
$r_h(d_h,S_h)$ and $r_f(d_f)$, so the government's cash cost of placing
$d_h$ and $d_f$ is the interest bill $r_h\,d_h + r_f\,d_f$. The
non-bank schedule depends on a country-specific structural size
$S_h>0$, which captures the depth of the domestic non-bank investor
base. The banking sector is not described by an interest rate:
placing $d_b$ with the banking system carries a separate welfare cost
$\Psi(d_b)$ that proxies for inflation, monetary--fiscal trade-offs,
and financial-stability risk. The planner therefore chooses
$d_h,d_f,d_b\ge 0$ to minimise total economic cost subject to the
budget constraint:
\begin{equation}\label{eq:planner}
  \min_{d_h,d_f,d_b}\;
    r_h(d_h,S_h)\,d_h + r_f(d_f)\,d_f + \Psi(d_b)
  \quad\text{s.t.}\quad d_h + d_f + d_b = D.
\end{equation}

%-------------------------------------------------------------
\subsection{The three sectors}
%-------------------------------------------------------------

Marginal cost rises with the sector's degree of saturation. What
differs is \emph{how} saturation is measured.

\paragraph{Non-banks.}
Non-banks operate in a segmented domestic market. Their interest rate
depends on the portfolio share $d_h/S_h$, not on $d_h$ alone: the same
absolute amount of debt is harder to place in a country with a thinner
non-bank investor base. We adopt the schedule
\begin{equation}\label{eq:nonbank-rate}
  r_h(d_h,S_h) = \bar r_h + \gamma_h\,
    \frac{d_h/S_h}{\bar\omega - d_h/S_h},
  \qquad 0\le d_h/S_h<\bar\omega,
\end{equation}
with $\bar r_h,\gamma_h,\bar\omega>0$. The rate stays near $\bar r_h$
when the sector is far from its country-specific cap $\bar\omega S_h$
and diverges as $d_h\to\bar\omega S_h$. A larger $S_h$ shifts the
capacity outward, so the same $d_h$ leaves more headroom and is
cheaper at the margin. Cross-country differences in $S_h$ are the
empirical hook of the model and the structural object recovered by
Stage~1 of the two-stage estimator in
Section~\ref{sec:empirical}.

\paragraph{Foreigners.}
Foreign participation is disciplined by a no-default / incentive
constraint that imposes a uniform upper bound $\bar D_f$ on foreign
holdings, common across countries. Risk premia rise as the sovereign
approaches that ceiling. We posit
\begin{equation}\label{eq:foreign-rate}
  r_f(d_f) = \bar r_f + \gamma_f\,
    \frac{d_f}{\bar D_f - d_f},
  \qquad 0\le d_f<\bar D_f,
\end{equation}
with $\bar r_f,\gamma_f>0$. The schedule is flat when
$d_f\ll\bar D_f$ (the global market is deep) and diverges as
$d_f\to\bar D_f$.

\paragraph{Banking.}
The banking sector absorbs the residual $d_b=D-d_h-d_f$. Unlike
non-banks and foreigners, it has no hard ceiling --- the central bank
can in principle keep expanding its balance sheet --- but doing so is
distortionary. We summarise that distortion by an increasing, convex
welfare cost
\begin{equation}\label{eq:banking-cost}
  \Psi(d_b),\qquad \Psi'(d_b)>0,\qquad \Psi''(d_b)>0.
\end{equation}
The level $\Psi'(0)$ is the cost of placing the first dollar with
banks and is directly comparable to $\bar r_h$ and $\bar r_f$. For
the numerical illustrations below we use the simplest
affine-marginal-cost specification
\begin{equation}\label{eq:psi-affine}
  \Psi(d_b) = \bar r_b\,d_b + \tfrac{\psi}{2}\,d_b^{2},
  \qquad \Psi'(d_b)=\bar r_b+\psi d_b,\qquad \psi>0,
\end{equation}
with $\bar r_b,\psi>0$.

%-------------------------------------------------------------
\subsection{Equilibrium and closed-form solution}
%-------------------------------------------------------------

With Lagrange multiplier $\lambda$ on the budget constraint, the
interior first-order conditions equate marginal placement costs across
all three sectors:
\begin{equation}\label{eq:foc}
  \text{MC}_h(d_h,S_h) \;=\; \text{MC}_f(d_f) \;=\; \Psi'(d_b)
  \;=\; \lambda(D),
\end{equation}
where, from \eqref{eq:nonbank-rate} and \eqref{eq:foreign-rate},
\begin{align}
  \text{MC}_h(d_h,S_h)
    &= \bar r_h + \gamma_h\,
      \frac{(d_h/S_h)\!\bigl[2\bar\omega - d_h/S_h\bigr]}
            {\bigl[\bar\omega - d_h/S_h\bigr]^{2}},
      \label{eq:MCh}\\
  \text{MC}_f(d_f)
    &= \bar r_f + \gamma_f\,
      \frac{d_f\,(2\bar D_f - d_f)}{(\bar D_f - d_f)^{2}}.
      \label{eq:MCf}
\end{align}
If any sector were cheaper at the margin, the planner would shift
toward it.

The system \eqref{eq:budget}--\eqref{eq:foc} admits a closed-form
parametric solution in the shadow rate $\lambda$:
\begin{align}
  d_h^\ast(\lambda,S_h)
    &= S_h\,\bar\omega\!\left(1 - \frac{1}{\sqrt{1+(\lambda-\bar r_h)/\gamma_h}}\right),
      \label{eq:dhstar}\\
  d_f^\ast(\lambda)
    &= \bar D_f\!\left(1 - \frac{1}{\sqrt{1+(\lambda-\bar r_f)/\gamma_f}}\right),
      \label{eq:dfstar}\\
  d_b^\ast(\lambda) &= \frac{\lambda-\bar r_b}{\psi},
      \label{eq:dbstar}
\end{align}
with $D(\lambda)=d_h^\ast+d_f^\ast+d_b^\ast$ strictly increasing in
$\lambda$ on the admissible range. Inverting numerically yields the
equilibrium allocation as a function of $D$.

%-------------------------------------------------------------
\subsection{Comparative statics: who absorbs the next dollar?}
%-------------------------------------------------------------

The substantive content of the model rests on a single structural
property of the three placement-cost schedules: in each sector, the
marginal cost of placing additional debt is strictly increasing in
that sector's existing holdings. Differentiating
\eqref{eq:MCh}--\eqref{eq:MCf} and \eqref{eq:psi-affine} gives the
slopes of the marginal-cost schedules,
\begin{align}
  H_h \;\equiv\; \frac{\partial\text{MC}_h}{\partial d_h}
    &= \frac{2\gamma_h\bar\omega^{2}}
            {S_h\,[\bar\omega-d_h/S_h]^{3}}>0,
      \label{eq:Hh}\\
  H_f \;\equiv\; \frac{\partial\text{MC}_f}{\partial d_f}
    &= \frac{2\gamma_f\bar D_f^{\,2}}
            {[\bar D_f-d_f]^{3}}>0,
      \label{eq:Hf}\\
  H_b \;\equiv\; \frac{\partial\text{MC}_b}{\partial d_b} &= \psi>0,
      \label{eq:Hb}
\end{align}
with $H_h\to\infty$ as $d_h\to\bar\omega S_h$ and $H_f\to\infty$ as
$d_f\to\bar D_f$. The banking slope is constant; the other two
diverge as their respective sectors saturate.

Total differentiation of the FOC \eqref{eq:foc} together with the
budget constraint \eqref{eq:budget} delivers the marginal share of
new debt absorbed by each sector,
\begin{equation}\label{eq:shares}
  s_j \;\equiv\; \frac{\partial d_j^\ast}{\partial D}
  = \frac{1/H_j}{1/H_h + 1/H_f + 1/H_b},
  \qquad j\in\{h,f,b\},
\end{equation}
which sum to one by construction. The marginal share of each sector
is therefore inversely proportional to the local curvature of its
marginal cost: at low $D$ all three $H_j$ are bounded and the $s_j$
are roughly constant; as $d_h\to\bar\omega S_h$ and
$d_f\to\bar D_f$, $H_h$ and $H_f$ diverge, $s_h$ and $s_f$ collapse to
zero, and $s_b\to 1$.

%-------------------------------------------------------------
\subsection{Three predictions}
%-------------------------------------------------------------

The structural property $H_j>0$ delivers three substantive
comparative-static results, which we record as propositions.
Proofs are collected in Appendix~\ref{app:proofs}.

\begin{proposition}[Sectoral holdings rise with total debt]
\label{prop:monotone}
Under upward-sloping marginal cost in every sector ($H_j>0$),
$d_j^\ast$ is weakly increasing in $D$ for $j\in\{h,f,b\}$, and
strictly increasing whenever $d_j^\ast>0$.
\end{proposition}

\noindent\emph{Intuition.} If, when the government must place more
debt, some sector's holdings actually \emph{fell}, that sector's
marginal cost would drop. The FOC requires all active sectors to
share the same marginal cost, so the others would have to drop too
--- meaning every sector absorbs less even though total debt has
risen. That is impossible, so no sector can shrink as $D$ grows.

\begin{proposition}[Non-bank absorption rises with non-bank sector size]
\label{prop:Sh}
Fix total debt $D>0$ and suppose the equilibrium is interior for
non-banks with $d_h^\ast>0$. Then
$\partial d_h^\ast/\partial S_h>0$.
\end{proposition}

\noindent\emph{Intuition.} Non-bank marginal cost depends on the
portfolio share $d_h/S_h$, not on $d_h$ alone: the sector becomes
expensive only as $d_h/S_h$ approaches $\bar\omega$. A larger $S_h$
means that any given level of $d_h$ corresponds to a smaller share,
so marginal cost is lower; holding $D$ fixed, the planner tilts the
allocation toward non-banks when $S_h$ rises. This is the mechanism
behind the cross-country pecking order in
Figure~\ref{fig:pecking-order}, and it is the structural object
recovered by Stage~2 of the empirical estimator.

\begin{proposition}[Strict concavity in the symmetric calibration]
\label{prop:concave}
Suppose $\bar r_h=\bar r_f$, $\gamma_h=\gamma_f$, and the equilibrium
path is interior for both capped sectors. Then $d_h^\ast(D)$ and
$d_f^\ast(D)$ are strictly concave in $D$ on the interior of their
domains.
\end{proposition}

\noindent\emph{Intuition.} Each sector's holdings rise with total
debt, but ever more slowly: placement is cheap when a sector is far
from its ceiling and expensive as it approaches $\bar\omega S_h$
(non-banks) or $\bar D_f$ (foreigners), so marginal issuance tilts
toward whichever margin still has room. Both capped curves therefore
bend downward, even though the government continuously substitutes
across sectors as relative costs shift.

\paragraph{Banking takeover at high debt.}
The three propositions sharpen into a single qualitative prediction
for the high-$D$ region. As $d_h\to\bar\omega S_h$ and
$d_f\to\bar D_f$, $H_h$ and $H_f$ diverge while $H_b$ remains
constant; \eqref{eq:shares} then forces $s_h,s_f\to 0$ and $s_b\to 1$.
Because the non-bank ceiling $\bar\omega S_{h,i}$ varies across
countries while $\bar D_f$ is common, each country has its \emph{own}
takeover level
$\bar D^{\text{bt}}_i\approx\bar\omega S_{h,i}+\bar D_f$ at which
banking begins to absorb the bulk of new issuance. Once $D$ exceeds
$\max_i\bar D^{\text{bt}}_i$, every country has exhausted its capped
sectors and banking absorbs nearly every additional dollar of debt
regardless of $S_h$. Cutting the horizontal axis at $D\approx 100\%$
of GDP therefore gives a uniform picture across countries: this is
the theoretical counterpart of the pooled empirical kink we estimate
in Section~\ref{sec:empirical}.

%-------------------------------------------------------------
\subsection{Numerical illustration: three phases of absorption}
%-------------------------------------------------------------

We illustrate the equilibrium with the baseline calibration
$\bar r_h=\bar r_f=\bar r_b=1$, $\gamma_h=\gamma_f=0.02$,
$\bar\omega=0.30$, $\psi=0.04$, and $\bar D_f=55$.
Figure~\ref{fig:phases} traces the equilibrium for a single country
with $S_h=150$ and shades the three phases defined by banking's
marginal share $s_b=\partial d_b^\ast/\partial D$. In Phase~1 both
capped sectors are far from their caps, $H_h$ and $H_f$ are bounded,
and \eqref{eq:shares} delivers a small $s_b$: non-banks and
foreigners absorb most new debt while banking is essentially flat.
As $D$ rises into Phase~2, $d_h$ and $d_f$ approach their caps, $H_h$
and $H_f$ rise sharply, and the planner shifts the marginal dollar
toward banking. In Phase~3 both capped sectors are close to their
limits, $1/H_h$ and $1/H_f$ collapse to zero, and $s_b$ approaches
one.

\begin{figure}[h!]
  \centering
  \includegraphics[width=0.85\linewidth]{\latexroot/Figures/absorption_phases.png}
  \caption{Three phases of debt absorption.}
  \label{fig:phases}
\end{figure}

%-------------------------------------------------------------
\subsection{Numerical illustration: heterogeneous non-bank sector}
%-------------------------------------------------------------

Figure~\ref{fig:pecking-order} overlays equilibrium allocations for
five countries that differ only in non-bank capacity,
$S_h\in\{30,70,110,150,200\}$, holding $\bar D_f=55$ and all other
parameters at the baseline. Because $1/H_h$ in \eqref{eq:shares} is
increasing in $S_h$, the model implies a pecking order in $S_h$:
at every level of $D$, countries with smaller non-bank sectors hold
a higher share of debt with foreigners and with banks. The middle
panel shows that small-$S_h$ economies climb the foreign schedule
faster and hit the uniform ceiling $\bar D_f$ at lower $D$; the right
panel shows that they also load more onto banking, and their banking
curves lie uniformly above those of large-$S_h$ countries. Once $D$
is large enough that every country has crossed its own
$\bar D^{\text{bt}}_i$, the five banking curves become approximately
parallel with $s_b\to 1$, while non-bank and foreign curves are flat
at their respective caps --- the cross-country signature of the
banking-takeover prediction.

\begin{figure}[h!]
  \centering
  \includegraphics[width=0.95\linewidth]{\latexroot/Figures/pecking_order.png}
  \caption{Pecking order across countries with different non-bank sector sizes $S_h$.}
  \label{fig:pecking-order}
\end{figure}

% Smart bibliography: only included when compiled standalone with citations
\smartbib

\end{document}
